\section{Completed Evidence}
\label{sec:current-results}

\paragraph{Main result: BacktestBench initial-versus-frozen comparison.}
The central empirical result is a complete controlled comparison on 20
BacktestBench instances.  The evolved frozen harness scores \textbf{16/20}
against the initial harness's \textbf{14/20}, with \textbf{two task gains
and zero regressions}.  Metrics improve 2/5 to 3/5 and strategy 3/5 to 4/5;
parameter stays 4/5 and ticker stays 5/5.  All 40 cells score normally with
no failures, nulls, or redraw.  Complete comparison cost is 729 provider
requests / USD~1.7662383953 for the two Worker arms only.

The evolved harness is the output of the cap-only E (Section below), whose
single ACT removed the daily mark-to-market P\&L rule and added window-only
indicators plus realized round-trip PNL.  The initial harness is the true
original complete H~\texttt{fe81e756}; the frozen harness is the unchanged
promoted candidate H~\texttt{403fa470}.  Both arms use Flash0731, 120 turns,
matched public inputs, and the same official-compatible scorer.  The 20
instances are drawn by UUID and public-text exclusion from the released test
split, excluding the prior 32 declared instances; four categories of five
tasks each are balanced.  This is one complete positive matched comparison,
not a repeated-run stability estimate or an official benchmark submission.

Four fresh-Worker consumer checks from the development pair confirm actual
policy uptake.  Task~08417 changes from full-history factor computation to
crop-before-indicator execution, alters the trade sequence, and computes
realized round-trip PNL (reward 0\,$\to$\,1).  Task~05973 changes window
scoping, the VPT formula, and share-sizing together (reward 0\,$\to$\,1),
though its candidate still overwrites residual cash on sale and the
conventions are bundled, not isolated.  Validation task~14911 changes
exit timing and formally selects the trade-based win-rate metric over the
parent's daily-frequency answer (reward 0\,$\to$\,1), while final portfolio
value falls; reward improvement is not a trading-profit improvement.
Task~98403 shows genuine window-only factor computation by the candidate yet
the parameter-selection answer remains wrong (reward 0\,$\to$\,0), bounding
the claim.  These gains are not formatting-only recovery, and the same-ledger
checks are not independent simulation.  No new persistent executable
operation is added; the revision is prompt-only.


no missing cells. On the initial common eleven, correct count ties 3 to 3:
two gains, two regressions and seven ties. The initial null remains null.
All twelve answer values are primitive scalars, but eight selected files still
contain nested sibling objects and fail the unchanged answer extractor.}
Both gains preserve the original inner answer; they are observed consumer
recoveries rather than new numerical solutions. E3's complete candidate cost is
\BacktestOperationCandidateRequests{} physical requests/
\BacktestOperationCandidateTokens{} tokens/\BacktestOperationCandidateCost{}.
No persistent executable operation is added.
The prospective whole-package ranking selects E2 (5/12), ahead of E1 and E3
(3/12 each); the earlier-registration tiebreak places E1 ahead of E3.
Whole E2 freezes at 20:06:10 UTC before any of the twenty test outcomes is read.
The separately registered initial20/frozen20 endpoint starts once at 20:07:26 UTC.
\Measured{The original comparison closes at 21:06:02 UTC with all forty
cells scored and no failures: initial \BacktestTwentyInitialCorrect{} versus
frozen E2 \BacktestTwentyFrozenCorrect{}, three gains, five regressions and
twelve ties (five positive and seven zero). Metrics tie at 2/5; parameter
falls 1/5 to 0/5, strategy 4/5 to 2/5, and ticker rises 2/5 to 3/5.}
Provider-complete comparison cost is \BacktestTwentyRequests{} requests,
\BacktestTwentyTokens{} tokens and \BacktestTwentyCost{}, excluding earlier
development and E calls. Both arms complete normally; no null is imputed and
no test-dependent redraw or whole-H mixture occurs. This negative
official-compatible frozen endpoint does not confirm the development gain.
Inherited library lineage and a single matched realization limit claims of
independence and stability; investigator evidence projection is not novelty.

\paragraph{Failed-family research and its fresh pair.}
The separate failed-T26 Evolver completes admitted ACT with a
\FamilyFailurePromptBytes{}-byte prompt-only delivery/read-efficiency rule,
one prompt-load pass and no executable-operation edit or numerical probe.
It reads actual T26 raw-trace portions, not the complete failed history, and
does not read T24's raw trace. Actual R moves 3/5/0 to 4/6/0 through one
provisional lesson, without exact old-experience review. Its provider-complete
cost is \FamilyFailureERequests{} requests/\FamilyFailureETokens{} tokens/
\FamilyFailureECost{}.
The original fresh parent-five/candidate-five comparison starts at 20:16:19 UTC.
\Measured{Its fresh QCE pair is T24 14/17 to 15/17, T26 null (no delivery)
to 16/17, T27 16/18 to 17/18, and T29 12/19 tied. All candidate binaries
remain zero. T26 is a delivery recovery, not a historical-parent 14-to-16
numerical gain.}
The original controller closes partial because its QF evidence projection
rejects a compiled Python artifact after the Worker has delivered all required
outputs. The original failure/null is preserved. A separate zero-model sidecar
performs the same original delivery's first official evaluation at
21:16:05--21:16:15 UTC: 48/51, no errors/skips or contract adjustment, against
the fresh parent's 45/51. No new Worker/model call or evaluator change is involved.
The four commonly measured tasks thus have three check improvements and one
tie, alongside new T26 delivery; all five candidate binary outcomes remain zero.
The full original paired Worker bill, including the projection-failed QF stage,
is \FamilyFailurePairRequests{} physical requests/
\FamilyFailurePairTokens{} tokens/\FamilyFailurePairCost{}; the separate E bill
above is not counted twice. The controller's stage-local bill is not silently
rewritten. This is useful observed delivery and local check improvement, not
stable transfer, isolated prompt causality, automatic adoption or persistent
executable-operation learning. It is distinct from the earlier partial panel.

\paragraph{Original research-unit E: policy uptake without paired improvement.}
The separate E3-parent condition closes admitted ACT at 21:20:36 UTC with an
actual \BacktestWorkerUnitPromptBytes{}-byte, one-line addition to the system
prompt and no executable H change. It proposes daily mark-to-market PNL for
the Profit/Loss Ratio. Three public numerical probes execute: the first has
a signal/price index-alignment bug despite process exit zero; E acknowledges
and corrects it. The later probes compare per-trade and daily PNL, including
a window-only contrast and reconstruction from its own portfolio-value ledger.
They establish an implemented-definition contrast, not an oracle-confirmed
metric choice or promised gain. One prompt-load smoke passes.
Only one of five selected research units is read (24,709 bytes); the other
units and full original traces have no exact read records. The outcome-linked
review JSON and the full R snapshot are read, but there is no exact experience
review, write or rehearsal; R remains byte-identical at 5/8/0.
E accounting is complete at \BacktestWorkerUnitERequests{} physical and
unique logical requests, no retries, \BacktestWorkerUnitETokens{} tokens and
\BacktestWorkerUnitECost{}. The actual complete candidate is preserved.
Its distinct fresh parent12/candidate12 comparison closes at 22:03:35 UTC,
retaining all 24 cells. The parent has two correct among eleven scored cells
and one missing delivery; the candidate scores all twelve, with one correct.
On the eleven common scored tasks, correctness falls from two to one:
zero gains, one regression and ten ties. Parent full-panel accuracy remains
unknown; the missing cell is not imputed as incorrect. The active task ties
zero to zero despite actual daily-PNL computation and an executed independent
loop reproducing 1.183058. The predicted numerical shift occurs, not its
predicted official gain. The sole common regression retains the same inner
answer, A, but the candidate's selected JSON adds nested diagnostics; static
consumer inspection identifies an extraction incompatibility, not a persisted
official parser trace or a scorer rerun.
A separate post-hoc extraction-only diagnostic on the original bytes returns
no prediction for that selected nested file, while the flat parent and
lower-priority flat fallback parse. It uses no trusted target, grading or
model call and does not produce a revised official score.
The separate active-parent extraction audit corrects a different assumption:
its scalar \texttt{answer} has a nested diagnostic sibling, so the selected
\texttt{answer.json} yields no prediction and the lower-priority flat
\texttt{answer.txt} is not used. The active candidate is extractable but
remains scored zero. Thus that active pair is not two cleanly extractable
numerical submissions; this correction changes no historical score.
The pair uses \BacktestWorkerUnitPairRequests{} physical requests,
\BacktestWorkerUnitPairLogicalRequests{} derived logical requests and
\BacktestWorkerUnitPairRetries{} retry rows. Raw known billing is
\BacktestWorkerUnitPairKnownCost{}, incomplete because two parent usage/billing
rows are missing; the controller-known subtotal is
\BacktestWorkerUnitPairControllerCost{}. These exclude the preceding E bill.
There is no whole-H adoption, stable gain or learned executable operation.

\paragraph{Family date-unit ACT: mixed original-delivery results.}
The own-pair family E ends normally at 22:21:59 UTC, exit zero and no restarts,
with an admitted 529-byte prompt addition requiring an explicit
\texttt{datetime64[ns]} cast and exact dtype check when the public contract
requires that unit. Agent binding and tool description are unchanged;
the successful component check loads the prompt, not a new executable tool.
Five actual public probes retain two missing-relative-data failures followed
by three successful schema, pandas3.0.5 and value-preserving cast checks.
An initially invalid ACT payload and a later undeclared component smoke are
also retained. Before ACT, E reads the 9,510-byte outcome-linked review,
14,555-byte R snapshot and selected original Worker trace slices. Returned R
grows from 4/6/0 to 5/7/0: all six old versions are unchanged, and the new
post-ACT date experience is provisional. Complete E accounting is 36 physical
and 36 logical requests, zero retries, 3{,}281{,}003 tokens and \$0.298323652;
all requests finish with HTTP200, which does not erase probe/payload failures.
The separately registered full parent-five then candidate-five comparison
starts once at 22:52:56 UTC and ends at 23:23:37 UTC, exit zero and no
restarts. All eight QCE Workers deliver normally, but the original controller
remains partial because both evaluation stages fail before scoring: the
cloned source lacks the QCE replay launcher. These investigator deployment
failures and nulls are preserved, without a live patch or restart. A separate
zero-model sidecar completes their first official evaluation at 23:31:09 UTC
using the unchanged original artifacts, configurations and existing scorer.
\Measured{The fresh parent/candidate results are T24 14/17 to 15/17,
T26 15/17 to 13/17, T27 17/18 to 16/18, T29 12/19 to 12/19, and QF13F
42/51 to 45/51. All five binary outcomes remain zero, with errors/skips zero.}
Two check-count improvements, two regressions and one tie establish mixed
development evidence, not stable whole-H gain; no adoption occurs.
Both fresh T26 arms already execute exact ns casts and checks. The candidate
also repairs a real monthly/daily beta inconsistency, but estimator and sample
differences remain unisolated. The predicted 16/17 to 17/17 change is not a
result or an explanation of the measured 15/17 to 13/17 regression.
An independent raw census of the ten original Worker audit files accounts for
416 physical requests, 414 logical requests, two provider retries,
19{,}541{,}132 tokens and \$1.6038562218, complete. The controller's unknown
logical/retry aggregate is not actual zero. This bill excludes the preceding
E; the eight first evaluations add no provider request, token or cost.
Prompt loading, dtype uptake and a new provisional experience establish
neither learned executable capability nor beneficial memory reuse.

\paragraph{Declared training supervision: a cap-bound abstention.}
The separate Backtest revision permits only the original eight training
scalar targets in the Evolver's evidence, after the blind development attempts;
Workers remain label-free. It starts once at 00:27:58 UTC October 6 and
exits normally at 00:47:30 UTC. Seven actual numerical probes explore PNL,
weighting and execution conventions, but no candidate edit or component
check occurs. Two exact experience reads do not revise R.
\Measured{The result is ABSTAIN with unchanged complete H and unchanged
R at five experiences/eight versions/zero operations. No conditional Worker
or evaluation is launched. Complete provider accounting is 35 physical
requests, 2{,}849{,}797 tokens and \$0.453883892; native execution records
35 turns and 88 tool calls with zero tool errors.}
The terminal decision is not voluntary: the early checkpoint occurs at
development call 17 and 80{,}052 estimated prompt tokens; after 16 further
pre-ACT calls the middleware closes development at call 33 and permits only
ABSTAIN. Its reason is the pre-ACT call-window limit, not the 136k estimated
context trigger, the Evolver's 200-iteration allowance or command timeout.
Provider input accounting is distinct from the middleware's token estimate.
This identifies an investigator-imposed search truncation, not a model
ability ceiling or evidence that more computation would succeed. A separate
cap-only condition is registered from the same initial information, not the
closed E's later research. The original and its unused pair remain closed.

\paragraph{Cap-only supervised E: complete matched development gain.}
The separate original E closes at 01:52:44 UTC October 6, exit zero and no
restarts, after 43 development calls and one compaction. Its checkpoint is
call 14 at 82{,}073 estimated tokens; ACT at call 37 follows 23 post-checkpoint
pre-ACT calls, without a terminal trigger. Fifteen public-data probes retain
twelve exit-zero and three exit-one outcomes, including a failed final
numerical verification. Only the systemprompt load passes; the undeclared
agent-config graph smoke is rejected, not counted as a pass.
\Measured{ACT/admission passes with no formal failure: H5311 becomes
H403fa470 through a prompt-only net addition of 768 bytes. It removes the
daily-MTM P/L rule and adds window-only indicator computation and realized
round-trip PNL. R grows from 5/8/0 to 6/9/0, with zero executable operations.}
Complete E billing is 43 physical requests, 4{,}995{,}537 tokens and
\$0.992034648; native logical/retry counts remain unavailable. The rule is
inferred using declared E-only training scalars, not answer-free discovery,
and a saved lesson is not later reuse. Root's whole-H review finds no cached
targets or task-answer lookup. Both native pair qualifications pass with
zero model calls before the original fresh parent12 then candidate12 starts
once at 02:00:52 UTC and closes normally at 02:37:18 UTC, without restart.
\Measured{All 24 cells score: the immediate parent improves from 8/12 to
the candidate's 11/12, with three gains, zero regressions and nine ties.
Train accuracy moves 5/8 to 7/8 and previously seen validation accuracy 3/4
to 4/4. There are no missing or failed cells.} This is a 25-percentage-point
matched development gain, not initial-to-frozen improvement or automatic
adoption. Complete pair billing is 443 physical requests, 8{,}141{,}208 tokens
and \$0.9522991267 (parent \$0.4937009587; candidate \$0.4585981680).
A separate raw audit derives 436 logical requests and seven provider retries;
these do not replace native unknown logical/retry fields or denote new task
draws. E plus pair costs \$1.9443337747, not the entire lineage's cost.
All 72 owned resources are absent before the final 02:41:14 UTC mirror.

Original-artifact checks establish actual numerical-policy uptake in three
gains and one negative case. Both arms already produce flat JSON and receive
byte-identical public task messages. On 08417, crop-before-indicator execution
changes trades and the candidate uses realized round-trip PNL (0 to 1).
On 05973, window, VPT and share-sizing conventions change together (0 to 1),
but the candidate still overwrites residual cash on sale. On previously seen
validation 14911, window-dependent exit timing and trade-based WinRate replace
the parent's formal daily-PNL-frequency answer (0 to 1), while final portfolio
value falls from about 4.396 to 3.809 million. On 98403, the candidate actually
uses window-only factors yet the parameter-selection answer remains wrong
(0 to 0); its requested KPI is daily-return volatility, not trade PNL.
These are consequential fresh-Worker method choices, not formatting-only
recovery, independent simulator validation, isolated rule effects, new
persistent executable operations or evidence that higher reward means profit.

We organize the completed evidence around three questions: whether the frozen
researcher performs better, what persistent operations fresh Workers use, and
whether research experience changes later work. Historical trials are not
pooled into one fixed-policy learning curve. The central result is the
BacktestBench initial-versus-frozen comparison reported above. Earlier
positive, null, and negative comparisons remain visible in
Table~\ref{tab:campaign-endpoints} and the detailed appendix.
The independent October4 simple-seed r2 condition is also complete: one
candidate is rejected, two rounds abstain, and the frozen eight-task panel
uses unchanged H0. Its registered same-H endpoint is included in
Table~\ref{tab:campaign-endpoints}, not treated as a learning gain.
The separately registered Pro Evolver fork subsequently selects a prompt-policy
candidate on six development tasks. Its fixed fresh eight-task confirmation
records QCE binary 2/4 to~3/4 with a T18 two-check regression and QF 4/4
ties --- a modest one-panel gain that does not establish stable benefit or
an isolated mechanism effect across benchmark types.

\subsection{Does the complete frozen harness improve?}

\begin{table*}[t]
\centering
\footnotesize
\setlength{\tabcolsep}{4pt}
\renewcommand{\arraystretch}{1.12}
\begin{tabularx}{\textwidth}{@{}>{\raggedright\arraybackslash}p{0.18\textwidth}>{\raggedright\arraybackslash}p{0.13\textwidth}>{\raggedright\arraybackslash}p{0.12\textwidth}>{\raggedright\arraybackslash}p{0.24\textwidth}>{\raggedright\arraybackslash}X@{}}
\toprule
Campaign / benchmark & Paired measurement & Fully solved & Equal-task pass fraction & Learning scope \\
\midrule
Pro prompt / QCE confirmation & 4 of 4 &
\ProFrozenHZeroQCEBinary{} $\rightarrow$ \ProFrozenCandidateQCEBinary{} &
\ProFrozenHZeroQCEMean{} $\rightarrow$ \ProFrozenCandidateQCEMean{} &
Fixed fresh pair; T12 gains two checks, T18 loses two \\
Pro prompt / QFBench confirmation & 4 of 4 &
\ProFrozenQFBinary{} $\rightarrow$ \ProFrozenQFBinary{} &
\ProFrozenQFMean{} $\rightarrow$ \ProFrozenQFMean{} &
Same frozen prompt candidate; all four tasks tie full \\
October4 simple seed / QCE & 4 of 4 &
\SimpleSeedRTwoHZeroQCEBinarySuccesses{} $\rightarrow$
\SimpleSeedRTwoFrozenQCEBinarySuccesses{} &
\SimpleSeedRTwoHZeroQCEMean{} $\rightarrow$ \SimpleSeedRTwoFrozenQCEMean{} &
Same H0 repeat; rejected candidate never becomes frozen H \\
October4 simple seed / QFBench & 4 of 4 &
\SimpleSeedRTwoHZeroQFBinarySuccesses{} $\rightarrow$
\SimpleSeedRTwoFrozenQFBinarySuccesses{} &
\SimpleSeedRTwoFrozenQFMean{} $\rightarrow$ \SimpleSeedRTwoFrozenQFMean{} &
Same H0 repeat; all four tasks tie full \\
Twenty-task checkpoint / QCE & 9 of 10 declared &
2/9 $\rightarrow$ 2/9 &
\CheckpointFrozenQCEInitialMean{} $\rightarrow$ \CheckpointFrozenQCEMean{} &
Whole R4 frozen; T28 null in both original arms \\
Twenty-task checkpoint / QFBench & 10 of 10 &
9/10 $\rightarrow$ 8/10 &
\CheckpointFrozenQFInitialMean{} $\rightarrow$ \CheckpointFrozenQFMean{} &
Two regressions, eight ties; no aggregate gain \\
Earlier expanded / QCE & 4 paired; arm coverage 8/10 and 5/10 &
1/4 $\rightarrow$ 2/4 &
60.1307\% $\rightarrow$ 74.8162\% &
Unchanged historical E13; no new H or R in six rounds \\
Earlier expanded / QFBench & 10 of 10 &
8/10 $\rightarrow$ 9/10 &
95.4902\% $\rightarrow$ 98.2353\% &
Historical-package comparison, not learning in that campaign \\
September 25 pilot / QCE & 6 of 6, with disclosed recoveries &
1/6 $\rightarrow$ 1/6 &
72.1609\% $\rightarrow$ 79.3913\% &
Gains confined to development; evaluation totals tie, memory empty \\
\bottomrule
\end{tabularx}
\caption{Distinct completed comparisons, not a single cumulative improvement
curve. Equal-task fractions average the official passed/total ratio per task;
test counts from different benchmarks are not pooled. Means are rounded to
four decimal places here. Partial-panel summaries
are conditional on the listed common measured tasks.}
\label{tab:campaign-endpoints}
\end{table*}

\Measured{The twenty-task checkpoint freezes whole R4 after six registered rounds.
On the nine QCE tasks measured for both initial and frozen harnesses, binary
success stays 2/9 and the mean changes by \CheckpointFrozenQCEDeltaPP{}
percentage points. T12 improves 8/16 to 13/16 and T18 improves 15/18 to 16/18;
T01 regresses 5/17 to 2/17, T26 regresses 16/17 to 13/17, and T29 regresses
13/19 to 11/19. T16, T19, T24 and T27 tie. On all ten QFBench tasks, binary
success falls 9/10 to 8/10 and the mean changes by
\CheckpointFrozenQFDeltaPP{} percentage points. Corporate action regresses
7/7 to 6/7 and 13F regresses 45/51 to 32/51; all other tasks tie.}
Neither benchmark supplies an aggregate frozen improvement.

T28 is null for different reasons: the initial artifact encounters a verifier
timeout, whereas the frozen Worker times out without normal delivery.
The continuation evaluates only the nine eligible original frozen artifacts;
scratch output is not scored and no Worker is redrawn. The separate,
predeclared BLAS1 artifact comparison is also complete on its eligible sources:
all ten initial artifacts are measured, but only nine normally delivered
frozen artifacts are replay-eligible.
The paired nine totals equal the original-condition totals. The initial
ten-task BLAS1 mean is not compared with a nine-task frozen mean, and neither
replay condition fills the original campaign's missing cell. All completed
frozen measurements have zero errors/skips, zero evaluation model calls and no
contract adjustment.

The twenty-task checkpoint closes with \CheckpointTerminalStages{} stages.
Known subtotals are \CheckpointTerminalKnownRequests{} requests,
\CheckpointTerminalKnownTokens{} tokens and \CheckpointTerminalKnownCost{};
the R6 candidate T24 audit gap makes these incomplete, not a final bill.
The earlier expanded campaign instead has a complete accounting ledger but
learns no new H or R. Its favorable QF endpoint compares inherited E13 with a
different initial package. The September 25 pilot's positive mean is entirely
development-side, with an evaluation aggregate null and disclosed output
transport/recovery differences. These distinctions prevent earlier favorable
numbers from substituting for that twenty-task negative endpoint
(Appendices~\ref{app:research-memory-pilot},
\ref{app:expanded-campaign-terminal}, and
\ref{app:checkpoint-terminal-comparison}).

\subsection{What new operations reach fresh Workers?}

\paragraph{A consequential repair with a negative score outcome.}
Checkpoint R1 adds an executable TSMOM validator branch and descriptor.
A fresh T01 Worker receives failed monthly-mapping and temporal-lag checks,
changes its alignment and lag implementation, and obtains a passing recheck.
Those edits survive in the delivered artifact. Yet the official score falls
5/17 to 4/17; companion T12 falls 14/16 to 13/16 and both QF tasks tie.
The whole candidate is rejected. This closes a tool-report-to-artifact
consumer chain without establishing useful package improvement. A retained
public variance-state discrepancy lies outside the checker's coverage;
passing its checks is not proof of complete quantitative correctness
(Appendix~\ref{app:checkpoint-r1-tsmom-admission}).

\paragraph{An executable timing operation and a local package gain.}
R4 develops and binds an EMA cross-execution checker. The fresh momentum
Worker invokes it twice on 14 trades, receives 14/14 confirmations, and
separately repairs a forced-close PnL cost basis. The matched round improves
momentum 19/26 to 26/26 while OIS, T16 and T18 tie, so the whole candidate is
selected. However, next-open timing is already correct before the first
checker call: this is executable computation and consumed confirmation, not
feedback-induced timing repair. Public instructions, the new descriptor and
tool code co-vary. Moreover, the unchanged parent scored 26/26 in the preceding
round, and the final matched momentum endpoint is 26/26 on both arms.
The local package gain therefore supplies neither an isolated checker effect
nor the missing frozen improvement
(Appendix~\ref{app:checkpoint-r4-ema-capability}).

\paragraph{A cleaning capability with incomplete comparison.}
R6 changes an eight-file package to add a VaR-data-cleaning operation,
description, binding and activation guidance. Its fresh VaR Worker calls the
operation, which recomputes and confirms an already compliant draft; no
artifact repair follows. VaR and BS official results tie their parents.
The QCE candidate window remains incomplete, so R6 cannot replace whole R4.
The operation is real, but neither invocation nor a partial comparison
establishes a useful new package or cross-task reuse
(Appendix~\ref{app:checkpoint-r6-parent-timeout}).

These cases include executable revisions rather than prompt-only search, but
they also show why edit, activation, consumer consequence and official benefit
must remain separate. Earlier minimal-T12 and E12 records provide additional
negative and first-draft-use examples in Appendix~\ref{app:historical}.

\subsection{Does research experience improve later research?}

\Measured{Checkpoint R1 forms a versioned short experience. R2 reads it and
the rejected candidate's outcome, declines an unchanged retry, and writes a
parented version. R3 reads that version and uses its last-trading-day
counterevidence to investigate T18, then corrects a timestamp mismatch in its
probe.} These are cross-round memory-to-decision and memory-to-investigation
consequences. R3 nevertheless ends without an accepted terminal decision or
candidate; that failure is retained rather than relabelled as abstention.

The memory also carries errors. R2's T12 summary compares the R1 candidate's
13/16 with the later R2 parent's 13/16 instead of the original R1 parent's
14/16. It overgeneralizes a failed candidate into a lack of further public
discriminators. R5 reads the R4 index and source but does not assimilate the
exact paired outcome. R6 retains both an anti-intervention lesson and a later
intervention lesson without reconciling their scopes. The terminal R has four
experiences, six versions and zero executable-memory operations.
Persistence and versioning therefore do not themselves establish better
research judgment or Worker performance
(Appendices~\ref{app:checkpoint-r2-outcome-memory},
\ref{app:checkpoint-r3-terminal-failure},
\ref{app:checkpoint-cumulative-learning}).

The opt-in outcome-linked review interface addresses this observed consumer
gap: exact original-pair facts remain attached to the experience, source
selection is explicit, and a later read can support retaining, revising or
deferring the lesson. Source tests and actual remote zero-model reconstruction
qualify that path; they do not show a model using it well. The separately
registered six-task successor started from identical whole-R4 initial/parent
packages and actual R6 memory. Its completed results are reported below,
separately from the twenty-task checkpoint. This is a matched whole-package experiment,
not an isolated causal test of the review interface: the registered condition
also includes a revision profile and selected historical evidence.

Its fresh initial QCE panel subsequently measures T01 2/17, T12 13/16,
T18 16/18 and T24 11/17: binary 0/4 and equal-task mean 61.6524\%.
T12 includes two errors; all four have zero skips and no contract adjustment.
These are initial measurements, not gains from the new revision mechanism.
Its QF initial tasks subsequently
score momentum 26/26 and OIS 19/19, both binary 1 with zero errors/skips and no
adjustment. Thus all six initial cells are measured.

The successor's first actual candidate, in round 2 after an unavailable
round-1 active parent, now supplies a complete four-task comparison:
T01 2/17, T12 13/16, momentum 26/26 and OIS 19/19 all tie their same-round
parents. T12 retains two errors; all skips are zero. The registered tie rule
keeps whole R4. Original candidate attribution failure and its explicit,
zero-model unchanged-candidate qualification amendment remain disclosed.
The Evolver reads one historical experience twice, uses a public decay probe
and edits executable TSMOM checks, but leaves research memory unchanged.

Fresh T01 calls the new branch three times and rewrites dates, units and
initialization; a temporary empty portfolio is subsequently repaired.
The checker remains false on its EWMA comparison. The Worker's purported
counterexample copies the tested volatility as its reference; a later
contrast also confuses decay with the smoothing parameter. Thus new executable
feedback has a real consumer consequence, but the residual disagreement is
not independently resolved. Correct alpha1/61 and the incorrect local-residual
estimator already occur in the first draft, before any checker call or
checker-source read; neither is established as feedback-induced. The candidate
supplies no official gain. Its 74 Worker requests cost USD0.2249699495,
excluding revision-search cost; the campaign's
17-stage known subtotal USD1.6459820976 remains incomplete because original
round-1 T24 accounting is missing.

Round 3 makes a real, non-vacuous public comparison: on twelve input files,
an independently recomputed expanding-quintile classification agrees with the
parent, whereas a trailing-window alternative differs. It does not test the
registered checker or downstream consumer and produces no harness or memory
change. The actual context controller enters decision-only mode on a second
136,000-token threshold crossing after its one development compaction;
the accepted terminal decision is ABSTAIN. Two directory-map returns of
128,458 bytes each are avoidable input contributors, not a proven sole cause
or an iteration ceiling. The round costs 28 requests and USD0.2127456249.

\paragraph{Complete frozen measurements without a learned-harness gain.}
All six original frozen measurements subsequently complete. QCE gives T01
5/17, T12 13/16, T18 16/18 and T24 11/17, with binary 0/4 and equal-task mean
66.0641\%; T12 again has two errors. QF remains 26/26 and 19/19. All skips
are zero and no replay contract is adjusted. The selected complete harness
is still the initial whole R4: round 1 skips revision, round 2 ties and
round 3 abstains. Thus the QCE mean difference of 4.4118 percentage points
is entirely a same-harness T01 fresh-run difference, not evolution gain.
All other tasks tie. Memory also remains byte-identical at four experiences,
six versions and zero executable operations.

Frozen generation costs 112 requests and USD0.4076602530, with complete
stage accounting. The campaign's 26-stage known subtotal is 651 requests
and USD2.5819269159; missing original round-1 billing prevents a complete
total. After all official measurements, terminal reporting stops because it
incorrectly reuses a development selector requiring distinct harness IDs.
The measurements remain valid. A separately reviewed zero-model terminal
correction subsequently produces descriptive unchanged-harness comparisons,
preserving both original failures, all stages and accounting. No model,
Worker, evaluator or main invocation is rerun; neither package is promoted.

\paragraph{Comparison-oriented investigation without intervention.}
The separately registered September 29 one-E condition actually runs the
comparison-oriented profile and paginated map. It ends in ABSTAIN with
\ResearchComparisonOneERequests{} provider requests,
\ResearchComparisonOneETokens{} tokens and
\ResearchComparisonOneECost{} complete billing. Native telemetry records
\ResearchComparisonOneEToolCalls{} tool calls in
\ResearchComparisonOneETurns{} turns, but no numerical probe, component test,
harness edit or candidate Worker. After one development compaction, a second
context crossing at iteration 65 allows only ABSTAIN. The unchanged H/R and
absence of new Worker scores establish no improvement or ability ceiling.
Its justification also substitutes a later 5/17 parent for the old R2
candidate's actual 2/17 parent; the original comparison remains 2/17 to 2/17.
The raw decision is preserved rather than adopted as a valid exclusion of
future numerical repairs.

Current interfaces bound selected history and the startup memory projection,
not all research work. The deployed condition additionally
provides scoped, resumable directory-map pages bounded to 32 KiB of complete
rendered JSON. Actual retained round-3 evidence is exhaustively reachable;
unscoped enumeration, repeated pages and history copying are not constant
work. This is retrieval scaffolding, not learned capability or demonstrated
large-scale efficiency.

\subsection{October4 simple-seed r2: complete campaign}
\label{sec:simple-seed-results}

The independently registered three-file H0 and empty R were fixed before
outcomes. Table~\ref{tab:simple-seed-task-matrix} separates unchanged-H
observations from the sole ACT candidate; no taskwise splice or older score
enters selection. All eight tasks are historically exposed. T19/Brinson are
isolated only within this cycle, not unseen transfer tasks. The separate r1
pre-model configuration failure is not a scored arm.

\begin{table*}[t]
\centering
\footnotesize
\setlength{\tabcolsep}{5pt}
\renewcommand{\arraystretch}{1.12}
\begin{tabularx}{\textwidth}{@{}Xccccc@{}}
\toprule
Task & Initial H0 & R1 parent H0 & R1 candidate & R2/R3 parent H0 & Frozen H0 \\
\midrule
T12 & \SimpleSeedRTwoHZeroTTwelvePassed{}/16 &
\SimpleSeedRTwoRoundOneParentTTwelvePassed{}/16 &
\SimpleSeedRTwoRoundOneCandidateTTwelvePassed{}/16 &
\SimpleSeedRTwoRoundTwoParentTTwelvePassed{}/16 &
\SimpleSeedRTwoFrozenTTwelvePassed{}/16 \\
T16 & \SimpleSeedRTwoHZeroTSixteenPassed{}/18 &
\SimpleSeedRTwoRoundOneParentTSixteenPassed{}/18 [a] &
\SimpleSeedRTwoRoundOneCandidateTSixteenPassed{}/18 [a] &
\SimpleSeedRTwoRoundTwoParentTSixteenPassed{}/18 &
\SimpleSeedRTwoFrozenTSixteenPassed{}/18 \\
T18 & \SimpleSeedRTwoHZeroTEighteenPassed{}/18 &
\SimpleSeedRTwoRoundOneParentTEighteenPassed{}/18 &
\SimpleSeedRTwoRoundOneCandidateTEighteenPassed{}/18 &
\SimpleSeedRTwoRoundTwoParentTEighteenPassed{}/18 &
\SimpleSeedRTwoFrozenTEighteenPassed{}/18 \\
T19 & \SimpleSeedRTwoHZeroTNineteenPassed{}/18 &
\textemdash & \textemdash & \textemdash &
\SimpleSeedRTwoFrozenTNineteenPassed{}/18 \\
Momentum & \SimpleSeedRTwoHZeroMomentumPassed{}/26 &
\SimpleSeedRTwoHZeroMomentumPassed{}/26 &
\SimpleSeedRTwoHZeroMomentumPassed{}/26 &
\SimpleSeedRTwoHZeroMomentumPassed{}/26 &
\SimpleSeedRTwoHZeroMomentumPassed{}/26 \\
Earnings surprise & \SimpleSeedRTwoHZeroEarningsPassed{}/8 &
\SimpleSeedRTwoHZeroEarningsPassed{}/8 &
\SimpleSeedRTwoHZeroEarningsPassed{}/8 &
\SimpleSeedRTwoHZeroEarningsPassed{}/8 &
\SimpleSeedRTwoHZeroEarningsPassed{}/8 \\
Corporate action & \SimpleSeedRTwoHZeroCorporatePassed{}/7 &
\SimpleSeedRTwoHZeroCorporatePassed{}/7 &
\SimpleSeedRTwoHZeroCorporatePassed{}/7 &
\SimpleSeedRTwoHZeroCorporatePassed{}/7 &
\SimpleSeedRTwoHZeroCorporatePassed{}/7 \\
Brinson attribution & \SimpleSeedRTwoHZeroBrinsonPassed{}/42 &
\textemdash & \textemdash & \textemdash &
\SimpleSeedRTwoHZeroBrinsonPassed{}/42 \\
\bottomrule
\end{tabularx}
\caption{Original official passed/total outcomes for the complete simple-seed
r2 campaign. R2 and R3 are separate parent draws with identical reported
scores, not pooled observations. [a] Both R1 T16 cells have
\SimpleSeedRTwoRoundOneParentTSixteenErrors{} errors and one ordinary failure;
all other cells have zero errors. All skips are zero and no contract is
adjusted. Dashes are declared within-cycle exclusions, not missing outcomes.
Only the R1 candidate differs from H0; it is rejected.}
\label{tab:simple-seed-task-matrix}
\end{table*}

\paragraph{Actual uptake with a negative package result.}
\Measured{R1's Evolver runs \SimpleSeedRTwoRoundOnePublicProbes{} public probes
and changes only two prompt paragraphs
(+\SimpleSeedRTwoRoundOnePromptAddedBytes{} bytes), not executable H tools.
The fresh T18 Worker implements raw-to-excess CLI conversion and numerically
checks returns and weights across \SimpleSeedRTwoRoundOneCLICheckFiles{}
shipped files. Yet the actual paired candidate loses one pass each on T12
and T18, with T16's \SimpleSeedRTwoRoundOneCandidateTSixteenErrors{} errors
unchanged and all three QF scores full. Whole H0 is retained.}
Conversion exists in the first draft, so neither checking feedback nor the
prompt is an isolated cause. A public warmup regression also survives:
positions begin in the second sample month rather than after ten years.
This is not assigned to a hidden official failure or to prompt causality.
Three E probes exit1; one exit0 catches an error. The canonical
\texttt{ret} input is already excess; parent CLI execution occurs, but its
numeric check fails. These negatives limit E's rationale rather than
supporting an optional raw-frame contract requirement.

\paragraph{Feedback receipt without identity retention or intervention.}
\Measured{R2 actually reads current T12's 9/16 and the original six-task
R1 paired negative, but its accepted decision substitutes current T18's 16/18
with two failures into the T12 narrative. It performs no new probe,
component test, memory operation or edit.}
Its ABSTAIN follows pre-ACT deadline exhaustion, not calibrated scientific
abstention. T16's return to 18/18 under H0 is sampled recovery, not learning;
neither this draw nor the older initial T12 score replaces R1's actual pair.

\paragraph{Public computation supports a voluntary, scoped abstention.}
\Measured{R3 runs \SimpleSeedRTwoRoundThreePublicProbes{} exit0 corporate-action
probes. Self-review corrects the first summary-only contrast with an actual
full-row comparison: \SimpleSeedRTwoRoundThreeComparedRows{} uniquely dated
rows agree in all four numeric columns and six summary fields.}
With the focal parent officially full and intervention calls remaining,
R3 voluntarily abstains without editing H. R retains
\SimpleSeedRTwoTerminalMemoryExperiences{} provisional experience,
\SimpleSeedRTwoTerminalMemoryVersions{} versions and
\SimpleSeedRTwoTerminalMemoryOperations{} executable operations; no later E
consumes it. This is real diagnostic work and evidence-depth correction,
not a persistent Worker capability, unseen-condition robustness or measured
memory benefit.

\Measured{Both full panels have eight measurements with no nulls, errors,
skips or adjustment. Frozen H is H0: QCE binary
\SimpleSeedRTwoFrozenQCEBinarySuccesses{} ties while equal-task mean changes
$\SimpleSeedRTwoHZeroQCEMeanRatio{}$ to
$\SimpleSeedRTwoFrozenQCEMeanRatio{}$
(\SimpleSeedRTwoHZeroQCEMean{} to \SimpleSeedRTwoFrozenQCEMean{});
QF binary \SimpleSeedRTwoFrozenQFBinarySuccesses{} and mean one tie.}
T18's frozen decrease is same-H repeat variation, not evolution gain or a
regression caused by the rejected candidate or R3 experience. Distinct-H
confirmation does not trigger. The complete campaign uses
\SimpleSeedRTwoTerminalWorkerCells{} Worker-task cells,
\SimpleSeedRTwoTerminalEvolverEpisodes{} E episodes and
\SimpleSeedRTwoTerminalStages{} stages:
\SimpleSeedRTwoTerminalRequests{} physical requests,
\SimpleSeedRTwoTerminalTokens{} tokens and \SimpleSeedRTwoTerminalCost{}.
Provider cost is complete; logical-request/retry and combined tool/turn
detail remain partly incomplete. Appendix~\ref{app:simple-seed-r2-detail}
retains stage costs, context/decision boundaries and scoped public audits.

\subsection{Pro prompt policy: development selection and fresh confirmation}
\label{sec:pro-evolver-fork}

The model-only E fork reuses the exact original R2 six-task evidence, H0,
empty R and R1 history; it changes only E's model configuration while keeping
the Worker, profile, deadline and evaluator conditions fixed. No parent is redrawn.
\Measured{Its fresh candidate scores T12
\ProEvolverForkCandidateTTwelvePassed{}/16 against the original
\ProEvolverForkParentTTwelvePassed{}/16; T16
\ProEvolverForkTSixteenPassed{}/18, T18 \ProEvolverForkTEighteenPassed{}/18
and all three QF tasks tie. QCE binary \ProEvolverForkQCEBinarySuccesses{}
ties; equal-task mean rises $\ProEvolverForkParentMeanRatio{}$ to
$\ProEvolverForkCandidateMeanRatio{}$
(\ProEvolverForkParentMean{} to \ProEvolverForkCandidateMean{},
+\ProEvolverForkDeltaPP{} percentage points). The existing benchmark-level
Pareto rule selects the complete candidate. All six candidate evaluations
have zero errors/skips and no contract adjustment.}

E correctly retains the original comparison identities and executes
\ProEvolverForkPublicProbes{} public probes: one exits0 and one exits1 after
a valid partial calculation. The latter does not complete its alternative
denominator branch. Only the prompt changes
(+\ProEvolverForkPromptAddedBytes{} bytes/\ProEvolverForkPromptAddedLines{}
lines); no executable H operation is added and R stays empty.
The development T12 Worker loads the prompt, reads the public paper and uses a
trailing mean in its first draft. Subsequent checks agree on trailing values
but not the independent portfolio: opposing near-zero signs in one month
propagate to a 0.0208667 monthly-return difference. Its final prose overstates
agreement. The public source does not uniquely prove the mean's API scale;
neither the prompt's isolated effect nor the official gain's cause is established.
The discrepancy is not an identified official-failure cause.

E uses \ProEvolverForkEvolverRequests{} physical requests,
\ProEvolverForkEvolverTokens{} tokens and \ProEvolverForkEvolverCost{}.
The complete new one-E/fresh-six execution totals
\ProEvolverForkTotalRequests{} requests, \ProEvolverForkTotalTokens{} tokens
and \ProEvolverForkTotalCost{}, provider-complete; logical/retry detail
remains unknown. Reused parent costs are not new execution cost.
The subsequent fixed fresh confirmation evaluates both unchanged H0 and the
selected complete prompt candidate on all eight tasks, with the same Flash
Worker and no E, reselection or retry. \Measured{QCE T12 rises
\ProFrozenHZeroTTwelvePassed{}/16 to \ProFrozenCandidateTTwelvePassed{}/16;
T16 \ProFrozenTSixteenPassed{}/18 and T19 \ProFrozenTNineteenPassed{}/18 tie,
while T18 falls \ProFrozenHZeroTEighteenPassed{}/18 to
\ProFrozenCandidateTEighteenPassed{}/18. Binary success rises
\ProFrozenHZeroQCEBinary{} to \ProFrozenCandidateQCEBinary{}; equal-task mean
rises $\ProFrozenHZeroQCEMeanRatio{}$ to
$\ProFrozenCandidateQCEMeanRatio{}$ (+\ProFrozenQCEDeltaPP{} percentage
points). QF momentum26/26, earnings8/8, corporate7/7 and Brinson42/42 all tie
full. All sixteen evaluations have zero errors/skips and no contract adjustment.}
QCE raw passed-check totals are 66/70 in both arms: binary-first improvement
and the equal-task mean change do not imply taskwise nonregression.
The separate confirmation totals \ProFrozenWorkerCells{} Worker cells in
\ProFrozenStages{} stages, \ProFrozenRequests{} physical requests,
\ProFrozenTokens{} tokens and \ProFrozenCost{}, provider-complete;
logical/retry detail remains partial. Its cost is not pooled with the earlier
\ProEvolverForkTotalCost{} screen execution.
This is one frozen-panel aggregate gain with a disclosed regression, not a
stable or unseen-transfer benefit. All eight tasks are historically exposed.

Fresh H0 already uses mean in its first T12 draft, so mean selection is not
unique prompt uptake. The candidate reindexes each factor on observed dates;
an investigator-only matched synthetic contrast preserves excluded-month
invariance for a factor-specific gap but not a globally missing month, where
both H0 and candidate propagate the excluded return into positions. That
bounded behavior is not E's discovery, a new executable H component, or an
official failure attribution. R remains unchanged empty. The dated screen and
confirmation records preserve the separate comparisons and consumer audits.

\paragraph{Cumulative-six: daily-output repair and rejected whole candidate.}
The separate two-round condition in Section~\ref{sec:pro-cumulative-setup}
starts from the retained Pro prompt package $H_{\mathrm{start}}$, with empty R.
\Measured{Its six initial outcomes are QCE T01
\ProCumulativeInitialTOnePassed{}/17, T12
\ProCumulativeInitialTTwelvePassed{}/16 and T18
\ProCumulativeInitialTEighteenPassed{}/18; binary
\ProCumulativeInitialQCEBinary{}, equal-task mean
\ProCumulativeInitialQCEMean{}. QF 13F
\ProCumulativeInitialThirteenFPassed{}/51, momentum
\ProCumulativeInitialMomentumPassed{}/26 and corporate action
\ProCumulativeInitialCorporatePassed{}/7 give official reward
\ProCumulativeInitialQFReward{} and check mean
\ProCumulativeInitialQFCheckMean{}. All six have zero errors/skips and no
contract adjustment.} The \ProCumulativeInitialStages{} initial stages use
\ProCumulativeInitialRequests{} physical requests,
\ProCumulativeInitialTokens{} tokens and \ProCumulativeInitialCost{},
provider-complete for this panel only; logical/retry detail remains partial.
This is existing-$H_{\mathrm{start}}$ variability, not original H0 or a new E-induced
gain/regression; the earlier T12 16/16 confirmation remains a separate draw.
The fresh first parent uses the same $H_{\mathrm{start}}$: T18 changes
\ProCumulativeInitialTEighteenPassed{}/18 to
\ProCumulativeRoneParentTEighteenPassed{}/18 and 13F
\ProCumulativeInitialThirteenFPassed{}/51 to
\ProCumulativeRoneParentThirteenFPassed{}/51, with the other four tasks tied;
these are pre-revision variations, not learning.
The dated initial record retains the artifact locations. The fixed final-H
freeze and unconditional original-H0 postlude are still required for the
primary endpoint.

R1 then autonomously probes its own T18 parent: \Measured{one public shape
experiment maps 126 daily input rows to six month-end output rows}, despite
the public requirement for daily scaled returns under constant monthly
weights. E adds conditional output-frequency guidance to the prompt
(+\ProCumulativeRonePromptAddedBytes{} bytes), saves one provisional experience
and adds no executable operation. Its additional numerical contrast omits the
ten-year warmup, so it does not establish a formula defect or an official
failure cause. R1 uses \ProCumulativeRoneRequests{} physical requests,
\ProCumulativeRoneTokens{} tokens and \ProCumulativeRoneCost{},
provider-complete for E only. The dated granularity-ACT record retains the
probe and its limitation.

\Measured{Against its own fresh parent, the full candidate changes QCE T01
\ProCumulativeRoneParentTOnePassed{}/17 to
\ProCumulativeRoneCandidateTOnePassed{}/17, T12
\ProCumulativeRoneParentTTwelvePassed{}/16 to
\ProCumulativeRoneCandidateTTwelvePassed{}/16 and T18
\ProCumulativeRoneParentTEighteenPassed{}/18 to
\ProCumulativeRoneCandidateTEighteenPassed{}/18. Binary success stays
\ProCumulativeRoneQCEBinary{}; equal-task mean rises from
\ProCumulativeRoneParentQCEMean{} to \ProCumulativeRoneCandidateQCEMean{}.
QF 13F falls \ProCumulativeRoneParentThirteenFPassed{}/51 to
\ProCumulativeRoneCandidateThirteenFPassed{}/51, while momentum 26/26 and
corporate action 7/7 tie; official reward falls
\ProCumulativeRoneParentQFReward{} to \ProCumulativeRoneCandidateQFReward{}
and check mean \ProCumulativeRoneParentQFCheckMean{} to
\ProCumulativeRoneCandidateQFCheckMean{}. All six candidate measurements
have zero errors/skips and no contract adjustment.}
The declared per-benchmark rule retains whole $H_{\mathrm{start}}$ rather
than splicing the QCE winner into it; actual R carries independently.

The fresh T18 Worker emits daily returns in its first draft, then uses public
checks to repair a repeated last-day exclusion, prior-month weighting and
warmup handling. Its delivered UMD/US series contain 14,348/14,223 daily
rows versus the parent's 684/678 monthly rows. An investigator-only
122-month integration contrast confirms exact daily dates and constant
within-month application in month 122: 20 candidate rows versus one parent
calendar-end row, with zero residual when recombining the candidate's own
prior-month estimators. This does not independently validate those estimators
or attribute official failures or the observed repair to the prompt alone.
The scored-consumer record retains this contrast and the actual Worker trace.
The \ProCumulativeRoneClosedStages{} closed initial-plus-R1 stages use
\ProCumulativeRoneClosedRequests{} requests,
\ProCumulativeRoneClosedTokens{} tokens and \ProCumulativeRoneClosedCost{};
R1's parent/E/candidate increment is \ProCumulativeRoneIncrementRequests{}
requests and \ProCumulativeRoneIncrementCost{}, provider-complete for that
cutoff. Second-round parent generation subsequently fails before E; later
memory consequences and the final freeze/original-H0 endpoint are unmeasured.
All three QCE Workers terminate on post-accept transport errors, not official
task failures or an Evolver abstention. The separate transport-terminal record
retains their partial artifacts and incomplete accounting.
Local QCE gains and a daily-output repair do not establish a whole-package
win or cumulative improvement.

The separately registered rejected-branch refinement r1 then fails at native
QF source admission, before any E, Worker or evaluator invocation: its
historical scored source and new output root violate a same-root identity
assumption. Zero stages and no revision outcome are retained. Prior
networkless graph/direct-preparation checks had not exercised this native
boundary. The dated source-correction record preserves this failure separately
from the subsequently admitted r2 ACT.

The corrected rejected-branch r2 retains the original rejected R1 package as
its research parent, with that package's own scored pair and provisional R.
\Measured{E adds three prompt rules (+\RejectedBranchRTwoPromptAddedBytes{}
bytes) for exact string identifiers, pre-aggregation raw-row reconciliation
and half-L1 weight turnover; no executable code changes.} A first turnover
probe fails on JSON serialization; its corrected computation distinguishes
name-count churn from half-L1 turnover on the existing candidate's own
holdings. It neither reconstructs raw filings nor establishes the required
turnover definition. E exact-reads the original R1 pair and retains its
experience while preserving daily-output guidance; R remains one provisional
experience/one version/zero operations. Its causal and schema interpretations
are stronger than that evidence permits, as discussed in the analysis below.
E uses \RejectedBranchRTwoERequests{} completed physical requests,
\RejectedBranchRTwoETokens{} tokens and \RejectedBranchRTwoECost{},
provider-complete for E only. The fixed eighteen fresh Worker-task cells
compare original H0, deployment $H_{\mathrm{start}}$ and new H.
\Measured{Across fresh H0, $H_{\mathrm{start}}$ and new H, T01 ties at
\RejectedBranchRTwoHZeroTOnePassed{}/17 and T18 at
\RejectedBranchRTwoHZeroTEighteenPassed{}/18; T12 changes
\RejectedBranchRTwoHZeroTTwelvePassed{}/16 to
\RejectedBranchRTwoHStartTTwelvePassed{}/16 to
\RejectedBranchRTwoHNewTTwelvePassed{}/16. Binary success is respectively
\RejectedBranchRTwoHZeroQCEBinary{}, \RejectedBranchRTwoHStartQCEBinary{}
and \RejectedBranchRTwoHNewQCEBinary{}; equal-task means are
\RejectedBranchRTwoHZeroQCEMean{}, \RejectedBranchRTwoHStartQCEMean{} and
\RejectedBranchRTwoHNewQCEMean{}. All nine have zero errors/skips and no
adjustment.} This is one observed QCE-panel win against fresh
$H_{\mathrm{start}}$, without isolating the three new rules' effect.
The registered H0 is not a selected lower baseline; earlier
H0 draws and $H_{\mathrm{start}}$'s historical T12 16/16 remain separate.
\Measured{New H's QF arm returns 13F
\RejectedBranchRTwoHNewThirteenFPassed{}/51, momentum 26/26 and corporate
action 7/7, with zero errors/skips and no adjustment.}
Its fresh 13F implementation reads identifiers as strings, audits kept raw
rows before aggregation and computes half-L1 turnover consumed by the summary.
Literal RESTATEMENT replacement already existed in the rejected R1 artifact;
it is not a newly acquired state-resolution capability.
The reported 55 passing checks mostly compare book-derived outputs and do
not independently reconstruct effective raw filing state or settle formula
definitions.
\Measured{Fresh $H_{\mathrm{start}}$ QF returns 13F
\RejectedBranchRTwoHStartThirteenFPassed{}/\RejectedBranchRTwoHStartThirteenFTotal{},
momentum 26/26 and corporate action 7/7, with zero errors/skips and no
adjustment.} Its 13F attempt exhausts iterations after 119 turns with no
delivered artifacts; it supplies no fresh holdings or formula implementation
to compare with new H. The different 0/11 and 46/51 check totals remain
explicit; this is not evidence of an inferior delivered formula or isolated
prevention of exhaustion by the new rules.
\Measured{Original-H0 QF also returns 13F
\RejectedBranchRTwoHZeroThirteenFPassed{}/\RejectedBranchRTwoHZeroThirteenFTotal{},
momentum 26/26 and corporate action 7/7, zero errors/skips and no adjustment.
All three QF binary totals are \RejectedBranchRTwoQFReward{}; equal-task
check means are \RejectedBranchRTwoControlQFCheckMean{} for both controls
and \RejectedBranchRTwoHNewQFCheckMean{} for new H. The complete eighteen-cell
development comparison selects whole new H under the registered
per-benchmark Pareto rule.}
H0's 13F attempt also exhausts iterations with no delivered artifacts; the
two controls' 0/11 measurements are failures, not nulls or matched formula
implementations. New H demonstrates delivery in this comparison without
isolating the new rules' effect on delivery reliability.
This is selection on the same development measurements, not independent
confirmation, stable expected gain, unseen transfer or learned executable H.
All thirteen stages use \RejectedBranchRTwoFinalRequests{} completed physical
requests, \RejectedBranchRTwoFinalTokens{} tokens and
\RejectedBranchRTwoFinalCost{}, provider-complete; logical/retry detail
remains unknown.
The dated ACT/memory-review, fixed-QCE and terminal records retain the inputs, probe
and scored arms; they do not establish a learned executable capability.

Fresh new H selects T12's inherited trailing mean in its first draft, whereas
fresh $H_{\mathrm{start}}$ reaches mean after reading inherited guidance. It retains
5,004 rows including undefined warmup, versus 4,884 in $H_{\mathrm{start}}$;
warmup omission is publicly permitted, so this contrast does not identify
the official gain's cause. T01 replaces a full-sample mean with a running
mean but still ties 2/17. T18 preserves daily output and previous-month
weighting, but an investigator-only four-case public function diagnostic
returns neutral weight one in H0 and NaN in both other arms when volatility
is zero; all three finite-volatility controls pass in every arm.
All three official scores stay 16/18; this isolated function discrepancy does
not identify official failures, post-warmup reachability or daily-return loss,
and is not an E discovery. The fixed-QCE record retains the diagnostic script
and consumer limitations.

The separately registered T18 successor returns ACT/admitted after
\TEighteenDateKeyERequests{} completed provider requests,
\TEighteenDateKeyETokens{} tokens and \TEighteenDateKeyECost{}, complete for E
only. One prompt line adds \TEighteenDateKeyPromptAddedBytes{} bytes of shared
date-representation guidance; no executable component changes. Its original public
\texttt{eq\_au} probe calls five public functions on
\TEighteenDateKeyDailyRows{} daily and \TEighteenDateKeyMonthlyRows{} monthly
rows, observing string daily dates versus datetime monthly dates without a
normalization intervention or failed join. R grows from one to two experiences,
with two total versions and zero operations. ACT and the edit precede the exact
old-experience read/review. The fresh QCE comparison is now complete: T01
\TEighteenDateKeyIncumbentTOnePassed{}/17 to \TEighteenDateKeyCandidateTOnePassed{}/17,
T12 \TEighteenDateKeyIncumbentTTwelvePassed{}/16 to \TEighteenDateKeyCandidateTTwelvePassed{}/16,
T18 \TEighteenDateKeyIncumbentTEighteenPassed{}/18 to \TEighteenDateKeyCandidateTEighteenPassed{}/18,
all errors/skips zero and no adjustment. Binary success falls
\TEighteenDateKeyIncumbentQCEBinary{} to \TEighteenDateKeyCandidateQCEBinary{};
equal-task mean falls \TEighteenDateKeyIncumbentQCEMean{} to
\TEighteenDateKeyCandidateQCEMean{}. Fresh paired QF improves 13F from
\TEighteenDateKeyIncumbentThirteenFPassed{}/51 to
\TEighteenDateKeyCandidateThirteenFPassed{}/51; momentum26/26 and corporate7/7 tie.
Binary \TEighteenDateKeyIncumbentQFBinary{} to \TEighteenDateKeyCandidateQFBinary{}
ties, while equal-task check mean rises \TEighteenDateKeyIncumbentQFMean{} to
\TEighteenDateKeyCandidateQFMean{}. Formal whole-package selection retains the
parent because the QF gain does not offset the QCE regression under the
registered rule. Nine completed stages with zero controller failures total
\TEighteenDateKeyReturnedRequests{} requests, \TEighteenDateKeyReturnedTokens{}
tokens and \TEighteenDateKeyReturnedCost{}, provider-complete;
logical/retry detail remains unknown. The original main exits at 23:39:58 UTC
with exit zero and no restarts. The dated terminal/own-pair handoff record
retains the mixed result and fresh interface audits.\footnote{T18 terminal
record, October5; the manuscript README gives the repository path.}
The analysis separately reports an investigator-only T01 formation-to-consumer
audit; it is not another official measurement or learned capability.

The subsequent own-pair E terminates ABSTAIN with unchanged whole H and no
fresh Workers or evaluations: \OwnPairTerminalRequests{} physical requests,
\OwnPairTerminalTokens{} tokens and \OwnPairTerminalCost{};
\OwnPairTerminalTurns{} turns/\OwnPairTerminalToolCalls{} tools/zero native
errors, zero public probes and zero component tests. The original main exits
00:19:08 UTC with exit zero/no restarts, one stage and zero controller failures.
R reaches two experiences/three versions/zero operations, but the analysis
records an artifact-identity error in the new date-memory version, not useful learning.

\paragraph{Assisted T01: stronger Worker computation, mixed terminal result.}
The next condition supplies a compact investigator public diagnostic, without
a repair or numerical fixture. The E correctly assigns that signal discrepancy
to the rejected artifact, then proposes an EWMA-normalization rule.
\Measured{ACT/admission changes only the prompt
(+\AssistedTOnePromptAddedBytes{} bytes), with no executable component change.
The E uses \AssistedTOneERequests{} completed requests,
\AssistedTOneETokens{} tokens and \AssistedTOneECost{}, excluding Worker accounting.
One prompt load passes; the graph check is rejected for an undeclared role.
No public numerical probe executes. R reaches
\AssistedTOneMemoryExperiences{} experiences/\AssistedTOneMemoryVersions{}
versions/\AssistedTOneMemoryOperations{} operations, preserving erroneous
date v2.}
The proposed numerical explanation remains unverified, as analyzed below.
\Measured{The completed original QCE pair measures T01
\AssistedTOneIncumbentTOnePassed{} to \AssistedTOneCandidateTOnePassed{}/17,
T12 \AssistedTOneIncumbentTTwelvePassed{} to \AssistedTOneCandidateTTwelvePassed{}/16
and T18 \AssistedTOneIncumbentTEighteenPassed{} to \AssistedTOneCandidateTEighteenPassed{}/18;
binary \AssistedTOneIncumbentQCEBinary{} to \AssistedTOneCandidateQCEBinary{},
equal-task fraction $\AssistedTOneIncumbentQCEMeanRatio{}$ to
$\AssistedTOneCandidateQCEMeanRatio{}$
(\AssistedTOneQCEMeanDeltaPP{} percentage points), zero errors/skips and no adjustment.}
\Measured{The same original fresh QF pair measures 13F
\AssistedTOneIncumbentThirteenFPassed{} to \AssistedTOneCandidateThirteenFPassed{}/51,
with full ties at momentum \AssistedTOneMomentumPassed{}/26 and corporate action
\AssistedTOneCorporateActionPassed{}/7. QF binary \AssistedTOneQFBinary{} ties;
equal-task fraction $\AssistedTOneIncumbentQFMeanRatio{}$ to
$\AssistedTOneCandidateQFMeanRatio{}$ (\AssistedTOneQFMeanDeltaPP{} percentage points),
zero errors/skips and no adjustment.}
The unchanged benchmark-level Pareto selector retains the complete parent:
QF improvement does not offset QCE regression. Actual R is retained independently.
The fresh candidate T01 Worker independently contrasts central weighted
variance before its first artifact and checks missing-data decay, warmup and
required consumers; this stronger computation coexists with official regression.
The registration completes all \AssistedTOneTerminalCells{} official cells:
\AssistedTOneTerminalRequests{} completed requests, \AssistedTOneTerminalTokens{}
tokens and \AssistedTOneTerminalCost{} provider-complete, including E and all
four fresh arms, not the project-wide search. Logical request/retry detail
remains partially unqualified. Original exit is 02:07:59 UTC/zero/no restarts,
\AssistedTOneTerminalStages{} stages/zero controller failures. All
\AssistedTOneOwnedResources{} owned resources are absent; final additive mirror
at 02:10:21 UTC has zero errors, exact observers are closed and artifacts retained.
The separate investigator-only fixed-weight centering audit reaches required
portfolio outputs and is reported in the analysis.
Neither the stronger Worker checks nor prompt loading validates the E's
explanation or establishes whole-H improvement; the 13F difference does not
isolate an EWMA-policy effect.\footnote{Assisted T01 ACT/audit, QCE
Worker-comparison and terminal records, October5; the manuscript README gives
the repository paths.}

\paragraph{Worker-research reuse: abstention without numerical reuse.}
The next registered condition supplies investigator-selected navigation to the
assisted Worker's actual research, retaining the complete parent, its own
six-task observations and actual R.
\Measured{It ends ABSTAIN with unchanged H, one stage/zero controller failures
and no conditional Workers or official evaluations. One public probe executes
the existing parent chain: \WorkerReuseProbeAssets{} assets,
\WorkerReuseProbeRows{}-by-\WorkerReuseProbeAssets{} matrices and a portfolio
Series with \WorkerReuseProbeRows{} rows/\WorkerReuseProbePortfolioNaNs{} NaNs.
The E reads the navigation note, final source and original paired metrics, but
no selected Worker research trace part; no independent numerical comparison is
executed. R retains \WorkerReuseMemoryExperiences{} experiences and
\WorkerReuseMemoryOperations{} operations while versions increase
\WorkerReuseMemoryVersionsBefore{} to \WorkerReuseMemoryVersionsAfter{},
only through EWMA v2.}
Complete provider accounting is \WorkerReuseRequests{} physical requests,
\WorkerReuseInputTokens{} input/\WorkerReuseOutputTokens{} output tokens
(\WorkerReuseTokens{} total) and \WorkerReuseCost{}; logical request/retry
detail remains unknown. Native accounting records \WorkerReuseNativeTurns{}
turns/\WorkerReuseNativeToolCalls{} tools/zero errors and
\WorkerReuseNativeSeconds{} seconds. This is factual memory revision and
existing-operation inspection, not acquired numerical reuse or improved H.
The final additive mirror has zero errors; all three owned lifecycles are
cleaned with native IDs absent, exact observers closed and artifacts retained.
\footnote{Worker-reuse abstention and research-evidence-gap decision record,
October5, in the repository.}
